\begin{examplebox}
	\textbf{\textcolor{CExample}{例1（基础）}}

	\textbf{\textcolor{CExample}{题目：}}
	已知直线$l$过点$(1,4)$，斜率$k=-2$，求$l$的方程。

	\textbf{\textcolor{CExample}{解题步骤：}}
	\begin{description}[style=nextline, leftmargin=2.8em, labelsep=0.5em, itemsep=0.45em]
		\item[\textbf{第一步（点斜式套用）：}] 由点斜式代入$x_0=1$，$y_0=4$，$k=-2$，得
			\[
				y-4 = -2(x-1).
			\]
			\par\textbf{理由：} 点斜式适用于已知一点和斜率。

		\item[\textbf{第二步（化为一般式）：}] 展开得$y-4=-2x+2$，移项得$2x+y-6=0$。
			\[
				2x+y-6 = 0.
			\]
			\par\textbf{理由：} 一般式是直线方程的规范形式，常作为最终答案。

		\item[\textbf{第三步（验证）：}] 将点$(1,4)$代入$2x+y-6$得$2\times1+4-6=0$，满足。
			\par\textbf{理由：} 确保方程正确。
	\end{description}

	\textbf{\textcolor{CExample}{关键结论：}}
	\textbf{$2x+y-6=0$}

	\medskip
	\textbf{\textcolor{CExample}{例2（稍变形）}}

	\textbf{\textcolor{CExample}{题目：}}
	已知直线$l$过点$(2,1)$和$(4,5)$，求$l$的方程。

	\textbf{\textcolor{CExample}{解题步骤：}}
	\begin{description}[style=nextline, leftmargin=2.8em, labelsep=0.5em, itemsep=0.45em]
		\item[\textbf{第一步（两点式套用）：}] 代入$x_1=2,y_1=1,x_2=4,y_2=5$，得
			\[
				\frac{y-1}{5-1} = \frac{x-2}{4-2} \;\Rightarrow\; \frac{y-1}{4} = \frac{x-2}{2}.
			\]
			\par\textbf{理由：} 两点式适用于已知两不同点。

		\item[\textbf{第二步（化为一般式）：}] 交叉相乘得$2(y-1)=4(x-2)$，展开得$2y-2=4x-8$，移项得$4x-2y-6=0$，化简得$2x-y-3=0$。
			\[
				2x-y-3 = 0.
			\]
			\par\textbf{理由：} 化为最简一般式。

		\item[\textbf{第三步（验证）：}] 将点$(2,1)$和$(4,5)$代入$2x-y-3$，满足。
			\par\textbf{理由：} 验证方程正确。
	\end{description}

	\textbf{\textcolor{CExample}{关键结论：}}
	\textbf{$2x-y-3=0$}

	\medskip
	\textbf{\textcolor{CExample}{例3（综合应用）}}

	\textbf{\textcolor{CExample}{题目：}}
	已知直线$l$过点$(1,2)$，且与直线$2x+3y-6=0$垂直，求$l$的方程。

	\textbf{\textcolor{CExample}{解题步骤：}}
	\begin{description}[style=nextline, leftmargin=2.8em, labelsep=0.5em, itemsep=0.45em]
		\item[\textbf{第一步（求已知直线斜率）：}] 将$2x+3y-6=0$化为斜截式: $y=-\frac{2}{3}x+2$，故斜率$k_1=-\frac{2}{3}$。
			\[
				k_1 = -\frac{2}{3}.
			\]
			\par\textbf{理由：} 一般式化为斜截式可得斜率。

		\item[\textbf{第二步（求垂直直线斜率）：}] 因为两直线垂直，斜率乘积为$-1$，所以$l$的斜率$k=\frac{3}{2}$。
			\[
				k = \frac{3}{2}.
			\]
			\par\textbf{理由：} 垂直条件$k_1\cdot k = -1$，解得$k$。

		\item[\textbf{第三步（点斜式写方程）：}] 由点斜式过$(1,2)$且斜率$\frac{3}{2}$，得$y-2 = \frac{3}{2}(x-1)$。两边乘$2$得$2y-4=3x-3$，移项得$3x-2y+1=0$。
			\[
				3x-2y+1 = 0.
			\]
			\par\textbf{理由：} 点斜式写方程，再化为一般式。
	\end{description}

	\textbf{\textcolor{CExample}{关键结论：}}
	\textbf{$3x-2y+1=0$}
\end{examplebox}
